% ECONOMICS_PROBLEM: {"id": "Q15-641", "title": "Land tenure and investment", "jel_code": "Q15", "source_id": "OP-0035"}
% !TeX program = xelatex
\documentclass[11pt,a4paper]{article}
\usepackage[margin=23mm]{geometry}
\usepackage{fontspec}
\setmainfont{Latin Modern Roman}
\usepackage{amsmath,amssymb,mathtools}
\usepackage{xcolor,enumitem,booktabs,longtable,array,xurl,hyperref}
\definecolor{muted}{HTML}{53636C}
\hypersetup{colorlinks=true,urlcolor=blue,linkcolor=blue}
\setlength{\parindent}{0pt}
\setlength{\parskip}{5pt}
\newcommand{\R}{\mathbb R}
\newcommand{\E}{\mathbb E}
\newcommand{\N}{\mathbb N}
\newcommand{\ind}{\mathbf 1}
\newcommand{\Var}{\operatorname{Var}}
\newcommand{\Cov}{\operatorname{Cov}}
\newcommand{\supp}{\operatorname{supp}}
\DeclareMathOperator*{\argmin}{arg\,min}
\DeclareMathOperator*{\argmax}{arg\,max}
\newcommand{\entrymeta}[3]{{\sffamily\small\color{muted}\textbf{\texttt{#1}}\quad|\quad #2\par #3\par}}
\newcommand{\Problemref}[1]{\texttt{#1}}
\newcommand{\JELref}[2]{\texttt{#2} (JEL \texttt{#1})}
\begin{document}
\section*{Land tenure and investment}
\textbf{Problem ID:} \texttt{Q15-641}\quad \textbf{JEL:} \texttt{Q15}\par
% BEGIN ECONOMICS STATEMENT
\label{problem:OP-0035}
\label{atlas:prob:D5}
\noindent\textbf{Original atlas ID:} D5 (entry 35). \textbf{Classification:} Editorial property-rights mechanism and equilibrium problem.

\noindent\textbf{Original atlas question:} When do stronger property rights raise investment, collateral use and productivity?

\subsection*{Definitions and assumptions}
Fix baseline parcels, occupants, households, and horizon $H$. Define a property-rights policy $r=(r_s,r_t,r_c,r_b)$ comprising tenure security, transferability, collateral eligibility, and intra-household or inheritance rights. The components need not change together. Let $I_{it}(r)$ be investment on baseline parcel $i$, $Q_{it}(r)$ physical or price-adjusted output, $B_{it}(r)$ borrowing under specified loan terms, and $A_{it}(r)$ occupancy or ownership status. Track displaced households as well as parcels. Potential outcomes may depend on the full community policy through land prices and insurance networks; specify interference assumptions or an equilibrium model $M$ including land and credit markets.
\subsection*{Formal research question}
For each feasible rights change $r_1-r_0$, identify or sharply bound
\[
 \Theta(r_1,r_0)=\left(
 \mathbb E\sum_{t=0}^H\beta^t[I_{it}(r_1)-I_{it}(r_0)],\;
 (\Delta Q_t,\Delta B_t,\Delta A_t)_{t=0}^H\right),
\]
where $\beta\in(0,1]$ and each expectation uses a stated baseline population. Determine how effects vary by initial tenure risk, lender enforceability, household bargaining, and land-market conditions. To separate channels, define component-specific interventions and measure their interactions; if only bundled title assignment is observed, report its bundle effect and bounds on the individual mechanisms rather than point estimates without identifying assumptions.
\subsection*{Interpretation and scope}
A legal title can change perceived security without becoming acceptable collateral; collateral eligibility can exist without profitable demand for borrowing. Loan availability is therefore distinct from borrowing and investment. Formal rights may crowd out informal insurance or change bargaining and concentration, so productivity gains need not imply gains for all occupants. Investment can itself increase de facto rights, creating reverse causality in observational comparisons. Titling assignment designs require credible counterfactuals and treatment versions, including enforcement and dispute resolution. Productivity estimates conditional on continued occupancy select on a policy outcome. Parcel outcomes, baseline-household welfare, and aggregate land-market effects should consequently be reported separately. A welfare comparison additionally requires costs of administration, displacement, and conflict, and declared social weights.
\subsection*{Source audit and status}
[S1] combines theory and Ghana evidence, [S2] studies urban Peruvian titling and labor supply, and [S3] studies a land-title natural experiment. Existing effects on these specific margins are not labeled unsolved. The component-specific and intergenerational equilibrium program is editorial and model-dependent. The sources do not imply that titles universally increase credit or investment, or certify a universal 2026 open frontier.

\subsection*{References}
\begin{enumerate}[label=\textnormal{[S\arabic*]},leftmargin=*]
\item Timothy Besley (1995). \emph{Property Rights and Investment Incentives: Theory and Evidence from Ghana}. Journal of Political Economy 103(5). \url{https://doi.org/10.1086/262008}. \textit{Locator:} Primary publisher, working-paper, or author record: bibliographic information and abstract; full-text theorem verification is not claimed. \textit{Role:} Investment/property-rights theory and Ghana evidence; emphasizes potential endogeneity of rights.
\item Erica Field (2007). \emph{Entitled to Work: Urban Property Rights and Labor Supply in Peru}. The Quarterly Journal of Economics 122(4), 1561--1602. \url{https://academic.oup.com/qje/article-abstract/122/4/1561/1850508}. \textit{Locator:} Primary publisher, working-paper, or author record: bibliographic information and abstract; full-text theorem verification is not claimed. \textit{Role:} Urban titling and labor-supply evidence; does not establish a universal collateral channel.
\item Sebastian Galiani and Ernesto Schargrodsky (2010). \emph{Property Rights for the Poor: Effects of Land Titling}. Journal of Public Economics 94(9--10). \url{https://doi.org/10.1016/j.jpubeco.2010.06.002}. \textit{Locator:} Primary publisher, working-paper, or author record: bibliographic information and abstract; full-text theorem verification is not claimed. \textit{Role:} Natural-experiment evidence on title assignment and household outcomes.
\end{enumerate}

\subsection*{Common conventions: Source mathematical conventions}

Unless an entry states otherwise, agent and firm sets are finite. State and action spaces are measurable, strategies and outcome maps are measurable, probability laws are specified, and expectations exist for the targets under discussion. Infinite-horizon discount factors lie in $(0,1)$ and discounted objectives are finite. Norms, tolerances and complexity input sizes must be specified for computational comparisons. Constants or primitives that appear locally are inputs, not empirical facts asserted by this catalogue.

\textbf{Identification.} A statistical model $m\in\mathcal M$ implies an observable law $P_m$ and target $\theta(m)$. At law $P$, its identified set is
\[
\Theta(P;\mathcal M)=\{\theta(m):m\in\mathcal M,\ P_m=P\}.
\]
Point identification holds when this set is a singleton, or equivalently when $P_{m_1}=P_{m_2}$ implies $\theta(m_1)=\theta(m_2)$ for all compatible models. Sharp bounds mean bounds attained, or approached when the boundary is not attained, by compatible models. Writing this set is a definition, not a solution: the substantive task is to establish informative restrictions and derive or estimate the set. An unrestricted-confounding model may give only trivial bounds.

\textbf{Inference.} Identification, estimability and valid uncertainty quantification are separate questions. Fix $\alpha\in(0,1)$. A confidence procedure $C_n$ over a declared law class $\mathcal P$ has asymptotic uniform coverage at level $1-\alpha$ when
\[
\liminf_{n\to\infty}\inf_{P\in\mathcal P}\Pr_P\{\theta(P)\in C_n\}\geq1-\alpha.
\]
For set-identified targets the coverage event must be defined separately, for example coverage of the full identified set. If an entry uses identification-indexed models instead of uniquely indexed laws, the same coverage convention is applied over those models. Pointwise limits do not establish uniform validity, and asymptotic validity does not guarantee finite-sample precision. Sample sizes, dependence structures and weak-identification sequences are part of each application.

\textbf{Counterfactuals.} $Y_i(a)$ is a potential outcome under a specified intervention, including its timing, implementation and versions. It is not $Y_i$ conditional on voluntarily choosing $a$. A full intervention vector permits interference; shorthand scalar treatment notation does not silently impose no interference. Assignment, selection, attrition, anticipation and measurement must be specified. A claim about an instrument's local effect is not automatically a population policy effect.

\textbf{Economic equilibrium.} A counterfactual model includes best responses, constraints, feasibility, market clearing, state transitions and expectations consistent with the stated information. When several equilibria exist, either a declared selector is an input or the target is set-valued. Comparisons across policy regimes must use the stated selection convention. Reduced-form relationships are not presumed invariant after an equilibrium-changing intervention.

\textbf{Optimization and welfare.} Optimization is over the stated feasible set. If attainment is not established, $\sup$ or $\inf$ and $\varepsilon$-optimal choices replace an asserted maximizer or minimizer. For example, with value $V=\sup_{a\in\mathcal A}W(a)<\infty$, an $\varepsilon$-optimal set is $\{a\in\mathcal A:W(a)\geq V-\varepsilon\}$ for $\varepsilon>0$. Benefits, costs, risk and health or environmental outcomes can be aggregated only using declared units and conversion weights. Interpersonal utility weights, the treatment of fiscal transfers and foreign profits, discounting, and the behavioral welfare criterion are explicit inputs. A comparative-static derivative additionally requires existence and appropriate differentiability of the selected counterfactual.

\textbf{Sources.} Local labels [S1], [S2], and so forth restart in every question. Locators identify the relevant abstract, model, section or result at the level actually checked; a bibliographic or abstract-level verification is not represented as inspection of an entire proof. Working papers are distinguished from later publications and changing versions. Source summaries are paraphrased. All URL checks and editorial formulations refer to the audit date stated above.
% END ECONOMICS STATEMENT
\end{document}
